\section{MARCO REFERENCIAL}

\subsection{Dominio de conocimiento seleccionado}

El dominio elegido para la práctica es el de los sistemas de alojamiento temporal de
intermediación tipo Airbnb. Este dominio involucra propiedades heterogéneas,
restricciones de uso, amenidades, perfiles de viajeros, finalidades de estadía y un
contexto geográfico que condiciona fuertemente la decisión de reserva. En consecuencia,
se trata de un escenario adecuado para aplicar modelado ontológico y recuperación de
información basada en semántica.

Una propiedad no puede analizarse únicamente por su precio o por su tipo. La decisión de
un usuario depende también de factores como la capacidad máxima, la cercanía a zonas
urbanas o puntos de interés, la disponibilidad de Wi-Fi, cocina o piscina, y la
compatibilidad con perfiles de viaje como pareja, familia, estudiante o ejecutivo. Esta
interdependencia entre atributos hace que el dominio sea más cercano a un problema de
representación de conocimiento que a un simple catálogo tabular.

\subsection{Justificación del dominio}

El dominio seleccionado es pertinente por cuatro razones principales:

\begin{itemize}
  \item \textbf{Complejidad conceptual}. El dominio incluye entidades distintas pero
  interrelacionadas: propiedades, amenidades, perfiles de huésped, propósitos de viaje,
  políticas y zonas geográficas.

  \item \textbf{Necesidad de consultas compuestas}. Las búsquedas reales combinan más de
  un criterio y requieren relacionar conceptos que no siempre se expresan como filtros
  exactos.

  \item \textbf{Disponibilidad de datos enlazados}. DBpedia provee recursos geográficos
  y etiquetas multilingües aprovechables para enriquecer el modelo local.

  \item \textbf{Valor demostrativo del multilingüismo}. El dominio es adecuado para
  mostrar internacionalización semántica, ya que las mismas entidades pueden presentarse
  correctamente en varios idiomas sin duplicar la ontología.
\end{itemize}

\subsection{Problema de los buscadores convencionales}

Los buscadores convencionales en plataformas de alojamiento suelen apoyarse en filtros
discretos y coincidencias literales. Este enfoque funciona para atributos explícitos,
pero falla cuando el usuario expresa necesidades semánticas. Por ejemplo, ``alojamiento
adecuado para trabajo remoto'' puede implicar Wi-Fi, espacio de trabajo, zona tranquila
o perfil ejecutivo; ``propiedad conveniente para estudiantes'' puede implicar precio
moderado, cercanía a universidades y estadías funcionales.

Cuando el modelo de datos no representa esas relaciones de forma explícita, el sistema
solo puede responder por coincidencias exactas o por reglas externas difíciles de
mantener. Una ontología permite, en cambio, representar el dominio como una red de
conceptos vinculados y consultables declarativamente mediante SPARQL.

\subsection{Conceptualización del dominio}

El análisis del dominio permitió identificar los conceptos centrales que estructuran el
modelo ontológico:

\begin{itemize}
  \item \textbf{Propiedad}. Representa la unidad ofertada, con atributos como nombre,
  descripción, precio por noche, capacidad máxima, calificación e imagen.

  \item \textbf{Tipo de propiedad}. Especializa el inventario en apartamento, casa,
  habitación privada, habitación compartida y villa.

  \item \textbf{Amenidad}. Modela servicios o equipamientos relevantes para la decisión
  de reserva, por ejemplo Wi-Fi, piscina, desayuno o espacio de trabajo.

  \item \textbf{Zona geográfica}. Describe la localización del alojamiento en una ciudad
  y una zona concreta, incorporando el contexto urbano del dominio.

  \item \textbf{Punto de interés}. Representa lugares próximos relevantes para la
  búsqueda, como plazas, universidades, terminales, hospitales o centros comerciales.

  \item \textbf{Perfil de huésped}. Describe el tipo de usuario o de grupo al que puede
  resultar más afín una propiedad.

  \item \textbf{Propósito de viaje}. Permite distinguir entre turismo, trabajo,
  estudios, congresos o reubicación.

  \item \textbf{Política}. Modela reglas de reserva o uso, necesarias para reflejar
  restricciones del dominio.
\end{itemize}

\subsection{Preguntas de competencia}

Las preguntas de competencia fueron definidas para asegurar que la ontología responda a
consultas representativas del problema. Entre las más relevantes se encuentran:

\begin{itemize}
  \item ¿Qué propiedades existen en una ciudad determinada y dentro de un rango de
  precio?
  \item ¿Qué alojamientos poseen amenidades específicas como Wi-Fi, piscina, cocina o
  espacio de trabajo?
  \item ¿Qué propiedades son compatibles con perfiles como familia, pareja, estudiantes
  o ejecutivos?
  \item ¿Qué alojamientos satisfacen simultáneamente restricciones de ubicación,
  capacidad, amenidades y contexto del viaje?
  \item ¿Cómo presentar la misma propiedad en distintos idiomas a partir de un único
  grafo RDF?
\end{itemize}

Estas preguntas guiaron tanto la estructura de la ontología como la implementación de
las consultas SPARQL y del prototipo web.

\subsection{Terminología base del dominio}

La terminología utilizada en el modelo fue refinada a partir del análisis del dominio,
de vocabularios de alojamiento y de las necesidades de recuperación expresadas por las
preguntas de competencia.

\begin{longtable}{|p{3.3cm}|p{5.2cm}|p{5.2cm}|}
\hline
\textbf{Área} & \textbf{Término en español} & \textbf{Equivalente operativo en inglés} \\
\hline
\endfirsthead
\hline
\textbf{Área} & \textbf{Término en español} & \textbf{Equivalente operativo en inglés} \\
\hline
\endhead
\hline
\multicolumn{3}{|r|}{\textit{Continúa en la siguiente página}} \\
\hline
\endfoot
\hline
\endlastfoot
Tipos de propiedad & Apartamento & Apartment \\
\hline
Tipos de propiedad & Casa & House \\
\hline
Tipos de propiedad & Habitación privada & Private room \\
\hline
Tipos de propiedad & Habitación compartida & Shared room \\
\hline
Tipos de propiedad & Villa & Villa \\
\hline
Amenidades & Wi-Fi & Wi-Fi \\
\hline
Amenidades & Cocina & Kitchen \\
\hline
Amenidades & Piscina & Swimming pool \\
\hline
Amenidades & Lavandería & Laundry \\
\hline
Amenidades & Desayuno & Breakfast \\
\hline
Amenidades & Espacio de trabajo & Workspace \\
\hline
Ubicación & Ciudad & City \\
\hline
Ubicación & Zona geográfica & Geographic area \\
\hline
Ubicación & Punto de interés & Point of interest \\
\hline
Perfil & Pareja & Couple \\
\hline
Perfil & Familia & Family \\
\hline
Perfil & Grupo estudiantil & Student group \\
\hline
Perfil & Ejecutivo & Executive traveler \\
\hline
Contexto & Viaje de trabajo & Business trip \\
\hline
Contexto & Viaje de estudios & Study trip \\
\hline
Contexto & Viaje vacacional & Vacation trip \\
\hline
Política & Política flexible & Flexible policy \\
\hline
Política & Política corporativa & Corporate policy \\
\hline
Política & Política para eventos & Events policy \\
\hline
Búsqueda & Consulta semántica & Semantic query \\
\hline
Búsqueda & Ranking de relevancia & Relevance ranking \\
\hline
Búsqueda & Filtros estructurados & Structured filters \\
\hline
\end{longtable}

\subsection{Relación con vocabularios y recursos existentes}

La ontología propia no fue construida aislada del ecosistema de Web Semántica. Se tomó
como referencia la existencia de vocabularios generales y específicos del dominio, así
como recursos disponibles en DBpedia. No obstante, dado que el objetivo de la práctica
era construir un modelo propio orientado a un buscador semántico concreto, se optó por
una ontología de dominio ajustada al caso de uso, y no por una simple reutilización
directa de un vocabulario externo completo.

Esta decisión permitió adaptar el modelo a necesidades reales del prototipo, como el
ranking semántico, los perfiles de huésped y la localización de zonas bolivianas, al
mismo tiempo que se preservó interoperabilidad mediante enlaces a recursos de DBpedia y
anotaciones multilingües estándar.

\subsection{Integración conceptual con DBpedia}

La integración con DBpedia se centró en los recursos más útiles y defendibles para el
dominio: tipos de alojamiento, amenidades y ciudades bolivianas. El criterio aplicado
fue seleccionar entidades de uso directo dentro del proyecto, evitando incorporar un
volumen excesivo de referencias irrelevantes para la demostración.

En términos prácticos, la integración se realizó mediante dos mecanismos:

\begin{itemize}
  \item \textbf{Alineación semántica}. Se añadieron enlaces
  \texttt{owl:equivalentClass} para clases como \texttt{Apartamento}, \texttt{Casa},
  \texttt{Villa} y \texttt{HabitacionPrivada}/\texttt{HabitacionCompartida} con sus
  equivalentes en DBpedia.

  \item \textbf{Enlace de individuos}. Se añadieron enlaces \texttt{owl:sameAs} para
  amenidades y ciudades bolivianas relevantes, aprovechando los recursos de DBpedia como
  referencia externa normalizada.
\end{itemize}

Este enfoque satisface el requisito de conexión con DBpedia y, al mismo tiempo, mejora
la robustez de la solución al mantener la operación principal sobre un dataset local en
Fuseki.

\subsection{Multilingüidad del dominio}

El enunciado exigía incorporar multilingüidad. En lugar de duplicar propiedades como
\texttt{nombreEn} o \texttt{descripcionEn}, se adoptó la estrategia recomendada en RDF:
etiquetar cada recurso con \texttt{rdfs:label} y \texttt{rdfs:comment} en varios
idiomas. Esto permite que el mismo individuo conserve identidad semántica única y sea
presentado en español, inglés o francés según la interfaz de usuario o la consulta
SPARQL ejecutada.

La elección del francés como tercer idioma, además de español e inglés, fortalece la
demostración de que la solución no fue diseñada solo para un caso bilingüe puntual, sino
para un esquema escalable de internacionalización semántica.
